Filtering news by the International Press Telecommunications Council (IPTC) Media Topic taxonomy needs a classifier that is cheap to run on a multilingual stream of articles. The open multilingual IPTC classifier, based on XLM-RoBERTa-large, takes \val{eff.teacher_onnx_int8.latency_ms_p50_1c} to \val{eff.teacher.latency_ms_p50_1c} milliseconds per article on one processor core. It is distilled into StaticKD, a student built on multilingual static embeddings whose linear head is folded into the embedding table, so the whole model is one table of 17 logits per token. The teacher labels \val{data.pool.docs} unlabeled articles in \val{data.pool.langs} languages. A test set of \val{ts.docs} articles in \val{ts.langs} languages, labeled by human annotators or by the sections that newsroom editors assigned and mapped onto the IPTC tree, allows evaluation without labels from any model. On this test set StaticKD reaches a macro-F1 of \val{statickd.test.macro}, or \val{statickd.test.pct_teacher} of the teacher score of \val{teacher.test.macro}. It classifies \val{eff.statickd.docs_per_sec_1c} articles per second per core, about \val{eff.statickd.speedup_onnx} times faster than the 8-bit quantized teacher, with a \val{final.size}-megabyte model, and it is more accurate than fastText and a linear support vector machine distilled on the same data. Ablations with five seeds per variant show that accuracy comes from multilingual distillation data and cross-lingually aligned embeddings, while hard labels match soft labels. Transfer to languages absent from the distillation data is carried by embedding rows that training never updates, and fine-tuning the embeddings trades this transfer for in-distribution fit. StaticKD is weakest on languages missing from both the distillation data and the pre-training data of the teacher.
